\section{Discussion} % (~20%)
% The discussion section is where you put your interpretations, extrapolations etc. It normally starts with concrete conclusions that come directly from your data itself, then gradually becomes more abstract as you extrapolate into the future and into broader research areas. You should start concrete, to show people that you are working with fact and not fantasy. But there is no harm in claiming potential greater significance for your results, if justifiable. Go back and check again, that what you have written in your discussion section ties in properly with the results you included in your report. What is the main message that you want readers to take home? Make this the primary focus of your discussion! The rest of the discussion should provide supporting arguments for this, your main thread. Your discussion section could address, e.g., select aspects of the robot's performance compared to the state-of-the-art and results and implications of your HRI experiment(s).

Now we can test our hypothesis H1: 

\par \textbf{Hypothesis H1:} Participants will more likely rate well the~robot's emotions that involve ear movements (happy, sad) compared to emotions that involve only head movements (interested, annoyed).


\subsection*{Method}
Ratings (Very Poor=1 to Very Good=5) were treated as ordinal data.  
Two independent groups:  
Group 1 (ear movements): Sad + Happy.  
Group 2 (no ear movements): Interested + Annoyed.  
Mann–Whitney \(U\) test, one-tailed (Group 1 > Group 2).

\subsection*{Result}
\[
U = 1122.0, \quad p \approx 0.187
\]
At \(\alpha = 0.05\), we fail to reject \(H_0\) (\(p > 0.05\)).

\subsection*{Conclusion}
Hypothesis H1 is **not supported** — there is no statistically significant evidence that emotions involving ear movements are rated higher than those with only head movements (\(p = 0.187\)).


\par 
